\section{Background and Related Work}
\label{sec:background}

\subsection{Existing judging platforms}

Commercial judging platforms exist (Judgify, Hackathon Junkies, Devpost) and
solve roughly the same problem we solve, but with different trade-offs. They
typically require an account for every voter, charge per event, and encode
their rubric as a fixed schema; which is fine when the rubric never
changes and the audience is captive but is the wrong shape for a one-off
academic event whose rubric depends on the segment type and whose audience
is fundamentally anonymous.

A second class of comparable artefacts are the open-source vote-collection
tools used in conferences (HotCRP for academic peer review, OpenReview for
ML conferences). These have very strong identity verification but no
real-time UI; they are batch-mode by design and therefore not applicable
to a live show.

\subsection{Real-time web stacks}

The traditional textbook real-time stack is WebSocket-based: Socket.IO on
top of Express, or a managed service like Pusher or Ably. WebSockets are
bidirectional and binary-friendly, which we do not need: every push in
Spotlight flows server-to-client. Server-Sent Events
(SSE)~\cite{w3c-sse} is the simpler primitive for one-way push, has wider
proxy support (it rides on plain HTTP/1.1 chunked encoding), and gives the
browser native auto-reconnect via the \texttt{EventSource} API. Where
SSE falls short; corporate proxies that buffer responses; a polling
fallback is sufficient because our update rate is bounded by human-event
cadence.

\subsection{Voting integrity literature}

The classic paper on rate-limiting and abuse detection in web systems
\cite{leiba-spam} long pre-dates modern bot frameworks but captures the
core idea: a layered defence is more cost-effective than a single perfect
check. Our anomaly detector follows that thinking: rate-limit at the
network edge (Express middleware), proof-of-presence at the protocol
edge (HMAC vote token), statistical scoring at the business layer
(burst detection on IP and subnet), and exclusion at the aggregate layer
(People's Choice ignores tainted projects).

CAPTCHAs are a closely related defence we considered but did not adopt
for the demo build. The argument against is twofold: (a) CAPTCHA UX
introduces friction in a 90-minute live event where every audience
interaction must be one tap, and (b) modern CAPTCHAs are bypassable
with cheap commercial solver services, which makes them more theatre
than barrier. The argument \emph{for} is that a CAPTCHA on
\texttt{/vote/init} would close the last avenue for a determined botnet,
and we acknowledge this gap explicitly in
Section~\ref{sec:future}.

\subsection{LLM-assisted evaluation}

Using an LLM as an advisory layer for human judges is a recent
practice; the consensus position in the academic-grading literature is that
LLM scores are unreliable as primary signal but useful as
\emph{question-prompts} or \emph{evidence-summarisers}~\cite{burstein-llm}.
Our AI Insights panel takes that position seriously: the panel never
returns a score. It returns a 60-word neutral summary, three apparent
strengths, three questions a judge could ask, and per-criterion hints.
The judge's score is the only score the system stores.

We use Claude Haiku 4.5 as the model. The prompt-caching
feature~\cite{anthropic-cache} lets us declare the rubric system prompt
\texttt{ephemeral}, which means the first call pays full input-token price
and subsequent calls within five minutes pay $\sim$10\% of that price for
the cached prefix. For an event with 9 featured projects and a recompute
on the day of the show, this turns nine paid-in-full calls into one paid
in full and eight effectively-free.

\subsection{Libraries we lean on}

\begin{itemize}[leftmargin=*]
  \item \textbf{React 19} with the \texttt{HashRouter} from React Router 7; chosen
        because Vercel's static hosting is happiest with a single
        \texttt{index.html} and we have no SEO requirement.
  \item \textbf{Vite 8} as the build tool. Module-level HMR makes the dev loop
        snappy and the production build emits a single
        \texttt{$\sim$740\,kB} bundle (gzipped to $\sim$217\,kB).
  \item \textbf{Express 4} for the API. Boring on purpose.
  \item \textbf{Supabase} as the Postgres host with a built-in REST/RPC layer.
        We use the service-role key from the backend; we never expose Supabase
        to the client.
  \item \textbf{Helmet 8}~\cite{helmet} for HTTP security headers, including a
        sensible default Content Security Policy.
  \item \textbf{express-rate-limit 7}~\cite{erl} for the per-route limits.
  \item \textbf{lru-cache 11}~\cite{lru-cache} for the in-process cache.
        Selected over \texttt{node-cache} because it has a stronger expiry
        story and a tighter API.
  \item \textbf{pino 9}~\cite{pino} for structured logging with
        \texttt{pino-http} for request-scoped log lines.
  \item \textbf{zod 3}~\cite{zod} for runtime validation. Acts as our
        compile-time check substitute.
  \item \textbf{@anthropic-ai/sdk 0.x} for the Claude API.
  \item \textbf{Recharts 2}~\cite{recharts} on the client for the analytics
        dashboard. Chosen over Chart.js because it composes naturally with
        React's render model.
\end{itemize}

\subsection{What is novel}

The novelty in Spotlight is not in any single component; the parts are
ordinary. The novelty is the \emph{combination}: a single small codebase
that simultaneously handles audience voting, judge grading, integrity
moderation, real-time admin push, and LLM-assisted reading aids, deployed
to free-tier infrastructure, and demonstrably defensible against the
non-functional probes a supervisor is likely to ask.
